\begin{center}
\textbf{\large{Tóm tắt}}
\end{center}

\addcontentsline{toc}{chapter}{Tóm tắt}

\begin{small}

Khóa luận này nghiên cứu bài toán chưng cất tri thức đa giáo viên không đồng
nhất cho mô hình thị giác-ngôn ngữ (VLM). Mục tiêu là truyền tri thức từ một
tập mô hình giáo viên đóng băng sang một mô hình học sinh nhỏ hơn khi các giáo
viên thuộc những họ mô hình khác nhau, sử dụng bộ tách từ khác nhau và có mức
độ hữu ích thay đổi theo từng mẫu. Trong bối cảnh này, chưng cất trực tiếp theo
chỉ số token không còn hợp lệ và phép lấy trung bình cố định giữa giáo viên có
thể trộn tín hiệu hữu ích với tín hiệu không liên quan hoặc xung đột.


Để giải quyết bài toán trên, khóa luận đề xuất \methodfull{}
(\methodname{}), một khung chưng cất cho truyền tri thức VLM đa giáo viên không
đồng nhất. \methodname{} kết hợp ba cơ chế: chưng cất Byte-Trie Wasserstein trên cây
byte UTF-8 dùng chung, định tuyến giáo viên thưa top-$k$ từ trạng thái ẩn của
mô hình học sinh và tinh chỉnh trọng số bằng độ đồng thuận gradient. Hàm mất
mát Byte-Trie Wasserstein giúp so sánh phân phối đầu ra của giáo viên và học sinh
giữa các bộ tách từ khác nhau. Định tuyến thưa chọn các giáo viên phù hợp với
từng mẫu. Tinh chỉnh bằng độ đồng thuận gradient sau đó điều chỉnh trọng số đã
định tuyến theo mức độ căn chỉnh giữa gradient chưng cất của từng giáo viên và
gradient cross-entropy có giám sát.


Thực nghiệm trên TextVQA, DocVQA và ChartQA đánh giá \methodname{} với hai mô
hình học sinh SmolVLM-500M và SmolVLM-256M trong thiết lập bốn giáo viên không
đồng nhất. Với SmolVLM-500M, \methodname{} đạt 69.38, 73.11 và 80.20 trên
TextVQA, DocVQA và ChartQA, so với điểm gốc 56.75, 59.52 và 62.80 của mô hình
học sinh. Các kết quả này tương ứng với mức cải thiện tương đối 22.26\%,
22.83\% và 27.71\% so với mô hình gốc, đồng thời thu hẹp 55.03\%, 39.53\% và
55.34\% khoảng cách ban đầu tới giáo viên có điểm cao nhất. Với SmolVLM-256M,
\methodname{} đạt 60.73, 60.48 và 72.32 trên cùng ba bộ đánh giá. Trên các
thực nghiệm được báo cáo, \methodname{} vượt qua tinh chỉnh trực tiếp, chưng
cất logits cùng họ mô hình, ULD khác họ mô hình và các phương pháp đa giáo
viên không đồng nhất được so sánh.

\vspace*{1cm}
\textbf{Từ khóa:} mô hình thị giác-ngôn ngữ, chưng cất tri thức, chưng cất đa
giáo viên, định tuyến giáo viên thưa, đồng thuận gradient, chưng cất
Wasserstein khác bộ tách từ

\end{small}
